% !PRINTED 468
\weberpage{468}

%% [running head, not body text -- commented out 2026-08-07] Page furniture (running
%% head, folio, signature) is excluded from this edition by standing policy.
%% \noindent 468 \hfill Nineteenth Chapter. \hfill \S\,127.

%<<CONT
by the one corresponding to the value $r=0$, or by the two corresponding to the values $r=0,2$, then it follows that $f(\omega)^6$ or $f(\omega)^{12}$ are class invariants.

If moreover $m$ is indivisible by 3, then one may assume the $B_r$ all divisible by 3, and it then follows, by combination with Theorem 1., that also
\[
f(\omega)^8f(\omega)^{-6}=f(\omega)^2 \quad\text{or}\quad f(\omega)^{12}f(\omega)^{-8}=f(\omega)^4
\]
are contained among the class invariants.

Now that our next goal has thus been indicated, we make use of the result of \S\,73, that between
\[
v=f\!\left(\frac{c+\partial\omega}{a}\right)^{3},\quad u=f(\omega)^3
\]
a modular equation:
\weq{7}{\Phi_n(u,v)=0}
holds, when
\[
a\partial=n,\quad c\equiv 0\pmod{16}
\]
The equation derived from this
\weq{8}{\Phi_n(u,u)=0}
is satisfied then and only then, when
\weq{9}{\dfrac{c+\partial\omega}{a}=\dfrac{\gamma+\delta\omega}{\alpha+\beta\omega},}
where $\begin{pmatrix}\alpha, & \beta\\ \gamma, & \delta\end{pmatrix}$ is a linear transformation of the first or second class which [according to \S\,40, (12)] satisfies the condition:
\[
\left(\frac{2}{\alpha-\beta}\right)e^{-\frac{\pi i}{8}(\alpha-\beta)(\alpha+\beta+\gamma-\delta)}=1,
\]
or the condition equivalent to it:
\weq{10}{\alpha\gamma+\beta\delta+2\alpha^2-2\alpha\beta-2\alpha\delta\equiv 0\pmod{16}.}
But from (9), if $\omega$ satisfies equation (2), there follows, as we have seen in the preceding section,
\weq{11}{
\begin{array}{ll}
n=y^2+mx^2, & \\
\alpha=Bx-y, & n\gamma=-Cx+c(Bx-y),\\
\beta=Ax, & n\delta=-Bx-y+Acx.
\end{array}
}
If we first assume $m$ to be odd, then we may set $n=m$, hence $x=1$, $y=0$, and since $c$ is divisible by 16, it follows:
\[
\alpha=B,\quad \beta=A,\quad \gamma\equiv -mC,\quad \delta\equiv -mB\pmod{8},
\]
%CONT>>
